% TOP_PROBLEM: {"id":30007651,"problem_number":"LOCAL-30007651","title":"Inheritance and long-run inequality","category_id":21,"category":{"id":21,"name":"economics","display_name":"Economics","description":"Open economic theory statements, published research agendas, and proposed scoped specifications; question provenance and historical priority are qualified per record.","slug":"economics","order_index":21,"created_at":"2026-10-08T00:00:00Z"},"status":"provenance_pending","external_url":null,"legacy_ids":[],"published":false,"statement":"In the explicitly specified setting: When do inheritances compress relative wealth inequality initially but increase long-run concentration through differences in heirs' saving behavior and investment returns? The mathematical statement above fixes the target, assumptions and scope of this catalogue formulation.","economics":{"original_id":140,"field":"Inequality and economic mobility","kind":"empirical","formalization_basis":"proposed_specification","source_status":"not_verified","priority_status":"not_established","priority_note":"No qualifying problem statement was directly verified, so historical priority cannot be assigned.","earliest_verified_reference":null,"current_reference":{"authors":["Salvatore Morelli","Brian Nolan","Juan C. Palomino","Philippe Van Kerm"],"title":"The influence of inheritances on wealth inequality in rich countries","year":2025,"url":"https://www.inet.ox.ac.uk/publications/the-influence-of-inheritances-on-wealth-inequality-in-rich-countries","doi":"10.1016/j.jpubeco.2025.105398","locator":"Abstract, paragraphs 1–3","evidence_summary":"The paper supplies size-dependent cross-sectional inheritance effects; the remaining long-horizon behavioral question is not explicitly stated here."},"formalization_note":"This is a proposed scoped research specification. No primary passage explicitly stating this entire remaining question as open was verified. The supporting literature may answer important subquestions; this entry must not be advertised as a verified formal open problem. The reference addresses cross-sectional size-dependent effects; it is not evidence for the earliest statement of the proposed dynamic mechanisms question. A corrigendum exists for the article."},"statusQualification":"Provenance pending: an explicit published statement of this exact open question was not verified. This is a catalogue-authored candidate specification, not a certified open conjecture. This is a proposed scoped research specification. No primary passage explicitly stating this entire remaining question as open was verified. The supporting literature may answer important subquestions; this entry must not be advertised as a verified formal open problem. The reference addresses cross-sectional size-dependent effects; it is not evidence for the earliest statement of the proposed dynamic mechanisms question. A corrigendum exists for the article.","statusReviewedAt":"2026-10-08"}
% FIRST_POSTED: 2026-10-08
% ENTRY_KIND: statement_only
% !TeX program = lualatex
\documentclass[11pt]{article}
\usepackage{fontspec}
\usepackage[a4paper,margin=25mm]{geometry}
\usepackage{amsmath,amssymb,mathtools}
\usepackage{xurl}
\usepackage[hidelinks]{hyperref}
\newcommand{\sref}[2]{\hyperlink{source:#1:#2}{[S#2]}}
\setlength{\emergencystretch}{3em}
\begin{document}
% problemId:30007651
\section*{Inheritance and long-run inequality}\label{30007651}
\subsection{Definitions and mathematical statement}
Fix an adult cohort with pre-inheritance wealth, parental links, and longitudinal portfolio records. Let \(b_i\) be person \(i\)'s net inheritance under a declared timing rule, and let \(b=(b_1,\ldots,b_N)\) be a feasible population-wide assignment for the \(N\) adults. Let \(W_{it}(b)\) be subsequent wealth allowing interactions among recipients. Its evolution includes labor income, consumption, realized portfolio returns, gifts, and mortality. For the population wealth distribution \(F_t(b)\), define \(\Delta_I(t,b)=I(F_t(b))-I(F_t(0))\), where \(I\) is a specified relative inequality measure. Estimate this path for different inheritance sizes, prior wealth ranks, and recipient ages. Determine how much variation in the path is accounted for by heirs' consumption and saving responses, business investment, portfolio choice, and returns, using additional identifying variation rather than accounting residuals alone. Preserve any established immediate equalizing or disequalizing effects in the relevant data. The research extension concerns dynamic behavioral mechanisms and their transportability; an observed cross-sectional marginal-transfer effect does not identify the long-run response. Death-related shocks may simultaneously alter grief, work, and transfers, so causal designs must address their exclusion restrictions. The cited 2025 paper answers size-dependent cross-sectional questions but an explicit remaining open-question passage for this dynamic formulation was not verified.

\par\textbf{Formulation and scope.} This is a proposed scoped research specification. No primary passage explicitly stating this entire remaining question as open was verified. The supporting literature may answer important subquestions; this entry must not be advertised as a verified formal open problem. The reference addresses cross-sectional size-dependent effects; it is not evidence for the earliest statement of the proposed dynamic mechanisms question. A corrigendum exists for the article.

\par\textbf{Open-question provenance.} An explicit published open-question statement was not verified. Sources below are supporting literature and must not be treated as a first listing.

\par\textbf{Historical priority.} No qualifying problem statement was directly verified, so historical priority cannot be assigned.

\subsection{Short English statement}
In the explicitly specified setting: When do inheritances compress relative wealth inequality initially but increase long-run concentration through differences in heirs' saving behavior and investment returns? The mathematical statement above fixes the target, assumptions and scope of this catalogue formulation.

\subsection{Sources}
\begin{itemize}\raggedright
\item[\textnormal{[S1]}]\hypertarget{source:30007651:1}{} Salvatore Morelli, Brian Nolan, Juan C. Palomino, Philippe Van Kerm. 2025. The influence of inheritances on wealth inequality in rich countries. Abstract, paragraphs 1–3.
\url{https://www.inet.ox.ac.uk/publications/the-influence-of-inheritances-on-wealth-inequality-in-rich-countries}.
DOI: 10.1016/j.jpubeco.2025.105398.
{\small\textit{Supporting or current literature; see evidence qualification. The paper supplies size-dependent cross-sectional inheritance effects; the remaining long-horizon behavioral question is not explicitly stated here.}}
\end{itemize}
\end{document}
